\documentclass[11pt]{article}
\usepackage[T1]{fontenc}
\usepackage{lmodern}
\usepackage[margin=1in]{geometry}
\usepackage{amsmath,amssymb,amsthm,booktabs,array,longtable}
\usepackage[hidelinks,breaklinks]{hyperref}
\newtheorem{theorem}{Theorem}
\newtheorem{lemma}[theorem]{Lemma}
\newtheorem{corollary}[theorem]{Corollary}
\newcommand{\ii}{\mathrm{i}}
\title{Extending Brown and Cameron's classification of purely imaginary independence roots}
\author{Haobo Ma\\\small ChronoAI Pte Ltd, Singapore
\and Wenlin Zhang\\\small National University of Singapore\\\small ChronoAI Pte Ltd, Singapore}
\date{November 19, 2026}
\begin{document}
\maketitle
\begin{abstract}
Extending Brown and Cameron's rational classification (2020,
arXiv:2001.08881v2, Proposition 2.9, p.~8), we answer their Question 2
(p.~11) affirmatively and classify all purely imaginary independence roots
whose imaginary part has rational square: they are exactly
$\pm\ii/\sqrt{n}$ for positive integers $n$.
An explicit connected graph of independence number four realizes each pair.
A second graph whose independence polynomial takes the value one at these
roots gives connected examples of every independence number at least four
and connected induced supergraphs of arbitrary finite graphs.
Bencs's tree-transfer theorem gives infinitely many trees for each pair.
A smaller 128-vertex example realizes $\pm\ii/\sqrt{2}$.
\end{abstract}

\section{The question and the classification}
Brown and Cameron ask whether a purely imaginary independence root can
have an irrational imaginary part \cite[Question 2, p.~11]{BC}.
We give an affirmative answer and determine every possibility with rational
squared imaginary part. Their Proposition 2.9 classifies the rational
imaginary parts as the nonzero reciprocals of integers. Its rational-root
argument already supplies the necessary restriction on rational squares;
the construction below realizes every permitted square, including all
nonsquare denominators.
The clique-substitution identity in Cameron's thesis
\cite[Theorem 3.3.4, p.~40]{Cameron}, together with its
$\pm\ii$ examples \cite[Proposition 3.4.1, p.~43]{Cameron}, already
implies the rational reciprocal examples; Brown and Cameron give their
explicit classification in 2020.

All graphs are finite and simple. For a graph $H$, an independent set is
a vertex subset containing no adjacent pair. Write $s_k(H)$ for the number
of independent $k$-subsets, $\alpha(H)$ for the largest such $k$, and
\[
 I(H,x)=\sum_{k=0}^{\alpha(H)}s_k(H)x^k.
\]
The empty set gives $s_0(H)=1$. An independence root is a complex zero of
this polynomial, and $\ii^2=-1$. The complete graph on $t$ vertices is
$K_t$; the edgeless graph on $t$ vertices is $E_t$ (including $E_0$).
For graphs on disjoint vertex sets, $H\sqcup J$ is their disjoint union,
and the join $H+J$ additionally joins every vertex of $H$ to every vertex
of $J$. Independent sets therefore give
\begin{equation}\label{eq:operations}
 I(H\sqcup J,x)=I(H,x)I(J,x),\qquad
 I(H+J,x)=I(H,x)+I(J,x)-1.
\end{equation}
In the second identity a nonempty independent set lies in exactly one
factor. These are the join and substitution methods used in \cite{BC}.

For each positive integer $n$, define the positive integers
\[
 r_n=81n^2-54n=27n(3n-2),\qquad d_n=27n^2+41n,
\]
and the graphs
\begin{align*}
 A_n&=K_{3n}\sqcup K_{3n}\sqcup K_{3n}\sqcup K_{3n},\\
 G_n&=A_n+(K_n\sqcup K_{r_n+1})+K_{d_n-1},\\
 J_n&=A_n+(K_n\sqcup K_{r_n})+K_{d_n}.
\end{align*}
Both joins have three nonempty factors; in particular $d_n-1\geq67$.

\begin{theorem}\label{thm:classification}
Let $b\neq0$ be real and suppose $b^2\in\mathbb Q$. Some finite graph
has independence root $\ii b$ if and only if $b^2=1/n$ for a positive
integer $n$. For every such $n$, the connected graph $G_n$ has
$108n^2$ vertices, independence number four, and
\begin{equation}\label{eq:general}
 I(G_n,x)=(1+nx^2)(1+108n^2x+81n^3x^2).
\end{equation}
The number $1/\sqrt n$ is irrational exactly when $n$ is not a perfect square.
\end{theorem}
\begin{proof}
For necessity, let
$E_H(t)=\sum_j s_{2j}(H)t^j$ be the even-coefficient polynomial of a
graph $H$. The real part of $I(H,\ii b)=0$ gives $E_H(-b^2)=0$.
This is an integer polynomial with constant term one. By the rational-root
theorem, the numerator of a rational root in lowest terms divides one.
Since $b\neq0$, its square is positive, so $b^2=1/n$ for an integer $n>0$.
This is the same necessity argument as \cite[Proposition 2.9]{BC}, without
requiring $b$ itself to be rational.

For sufficiency, use $I(K_t,x)=1+tx$ and \eqref{eq:operations} to obtain
\begin{align*}
 I(G_n,x)&=(1+3nx)^4+
       \bigl((1+nx)(1+(r_n+1)x)-1\bigr)+(d_n-1)x\\
 &=1+108n^2x+(n+81n^3)x^2+108n^3x^3+81n^4x^4.
\end{align*}
Here the linear coefficient is
$12n+n+r_n+d_n=108n^2$ and the quadratic coefficient is
$54n^2+n(r_n+1)=81n^3+n$. Multiplication proves
\eqref{eq:general}. At $z=\pm\ii/\sqrt n$, the factor $1+nz^2$ vanishes.

The order is $12n+n+(r_n+1)+(d_n-1)=108n^2$.
The final nonempty clique is adjacent to every vertex in the other
factors, so $G_n$ is connected. An independent set in $A_n$ has at most
one vertex from each of four cliques, and such a four-set exists.
The other two join factors have independence numbers two and one.
Thus $\alpha(G_n)=4$.

Finally, if $1/\sqrt n=u/v$ for coprime positive integers $u,v$, then
$nu^2=v^2$. A prime dividing $u$ would also divide $v$, so $u=1$ and
$n=v^2$. Conversely a perfect-square $n$ gives a rational reciprocal
square root. This proves the last assertion.
\end{proof}

\section{Connected completions and trees}
Theorem~\ref{thm:classification} in particular gives $|b|\leq1$ when
$b^2$ is rational. Cameron's separate question about $|b|>1$
\cite[fourth question]{CameronQ} is therefore not settled here:
any such example would have irrational $b^2$.

The companion graph $J_n$ permits connectedness without changing the roots.
\begin{lemma}\label{lem:neutral}
The graph $J_n$ has $108n^2$ vertices, independence number four, and
\begin{equation}\label{eq:neutral}
 I(J_n,x)=1+(1+nx^2)(108n^2x+81n^3x^2).
\end{equation}
Consequently, if $I(Q,z)=0$ and $z=\pm\ii/\sqrt n$, then $Q+J_n$
is connected and has independence root $z$ and independence number
$\max\{\alpha(Q),4\}$.
\end{lemma}
\begin{proof}
Relative to $G_n$, the factor $K_n\sqcup K_{r_n+1}$ loses one vertex
in its second clique, while the final clique gains one. The changes
to their independence polynomials are $-x-nx^2$ and $+x$.
Thus $I(J_n,x)=I(G_n,x)-nx^2$, which is \eqref{eq:neutral}.
The order is unchanged; the same join-factor argument gives independence
number four. At either specified $z$, equation \eqref{eq:neutral}
gives $I(J_n,z)=1$, so \eqref{eq:operations} gives
$I(Q+J_n,z)=I(Q,z)+1-1=0$.
A graph with such a root is nonempty, and $J_n$ is nonempty, so their
join is connected. Independent sets in a join stay in one factor,
giving the stated maximum.
\end{proof}

\begin{corollary}\label{cor:completion}
Fix $n\geq1$.
For every integer $a\geq4$ there is a connected graph of independence
number $a$ with roots $\pm\ii/\sqrt n$.
Every finite graph $H$ is an induced subgraph of a connected graph
with these same roots.
\end{corollary}
\begin{proof}
For the first assertion take
$(G_n\sqcup E_{a-4})+J_n$.
Before the join its polynomial is $I(G_n,x)(1+x)^{a-4}$ and its
independence number is $a$. Apply Lemma~\ref{lem:neutral}.
For the second take $(G_n\sqcup H)+J_n$.
The disjoint union retains the roots of $G_n$, and the join does not
alter adjacency among the vertices of $H$.
This graph has $216n^2+|V(H)|$ vertices and independence number
$4+\alpha(H)$, also when $H$ is empty.
\end{proof}

These extend the rational-root conclusions of
\cite[Theorems 2.8 and 2.12]{BC} to every rational squared imaginary part
permitted by Theorem~\ref{thm:classification}. The existence of independence
number four examples is consistent with the obstruction
$\alpha\leq3$ in \cite[Proposition 2.1]{BC}.

\begin{corollary}\label{cor:trees}
For each positive integer $n$ there are infinitely many pairwise
nonisomorphic finite trees with independence roots $\pm\ii/\sqrt n$.
None of these trees has a universal vertex.
\end{corollary}
\begin{proof}
Bencs \cite[Proposition 2.7(1)]{Bencs} proves that for every connected
finite graph $H$ there are a tree $T$ and induced subgraphs
$H_1,\ldots,H_t$ of $H$ such that
\[
 I(T,x)=I(H,x)\prod_{j=1}^t I(H_j,x).
\]
This is also \cite[Proposition 3.1]{BC}.
Apply it to the connected graphs of independence number $a$ in
Corollary~\ref{cor:completion}, for arbitrarily large $a$.
Their specified roots are roots of the tree polynomials. Since
$\deg I(T,x)=\alpha(T)\geq\alpha(H)=a$, the resulting trees have
unbounded independence number, hence infinitely many isomorphism types.
This extends \cite[Corollary 3.2]{BC}, which treats $\pm\ii/m$.
A tree with a universal vertex is a star, with polynomial $(1+x)^t+x$.
If $t>0$ and $z=\pm\ii/\sqrt n$, then $|1+z|^t>1\geq|z|$, so
the polynomial cannot vanish. For $t=0$ the graph is $K_1$, whose root
is $-1$. Thus any tree with one of the specified roots has no universal vertex.
\end{proof}

\section{A smaller irrational example}
Let $A$ be four disjoint copies of $K_4$, let $B$ be the join of
34 copies of $E_2$, and set $W=A+B+K_{44}$.
The pairs in $B$ are completely adjacent to one another.
Then \eqref{eq:operations} gives
\begin{align*}
 I(W,x)&=(1+4x)^4+34(2x+x^2)+44x\\
 &=1+128x+130x^2+256x^3+256x^4\\
 &=(1+2x^2)(1+128x+128x^2).
\end{align*}
Thus $W$ has roots $\pm\ii/\sqrt2$. Its order is $16+68+44=128$,
its final clique makes it connected, and its join factors show that its
independence number is four. This is smaller than $G_2$, which has 432
vertices. Clique substitution replaces each vertex of $H$ by $K_m$,
joining replacement cliques according to the edges of $H$. Every independent
$k$-set lifts in $m^k$ ways, so $I(H[K_m],x)=I(H,mx)$, the clique case of
Brown--Hickman--Nowakowski's substitution theorem \cite[Theorem 1]{BHN},
also \cite[Theorem 2.6]{BC}. In particular $W[K_m]$ has roots
$\pm\ii/(m\sqrt2)$ for every positive integer $m$.

\section{Source coverage and mathematical support}
Table~\ref{tab:coverage} records the related source statements.
The posing locator is Question 2 on printed p.~11 of arXiv v2, not a
journal-page locator. The table's extensions are proved above; prior
results retain their original attribution.
\small
\begin{longtable}{@{}p{0.42\linewidth}p{0.53\linewidth}@{}}
\caption{Coverage of the posing source and closely related statements.}\label{tab:coverage}\\
\toprule Source item & Coverage\\\midrule\endfirsthead
\toprule Source item & Coverage (continued)\\\midrule\endhead
\bottomrule\endfoot
Brown--Cameron Question 2, p.~11: irrational imaginary part
 & Answered: yes. Theorem~\ref{thm:classification} and the example $W$.\\[4pt]
Proposition 2.9, p.~8: rational imaginary parts; exclusion after Question 2
 & Reproved; prior proof Brown--Cameron (2020). Theorem~\ref{thm:classification}
 extends the rational-root restriction to a classification for $b^2\in\mathbb Q$.\\[4pt]
Theorems 2.3 and 2.5: connected examples, $\alpha=4$ and every $\alpha\geq4$
 & Prior proof cited for roots $\pm\ii$; Corollary~\ref{cor:completion}
 extends to $\pm\ii/\sqrt n$.\\[4pt]
Earlier v1, Proposition 2.1 and Question 1, pp.~3 and 9:
 exclusion boundary and possible independence numbers
 & Prior correction and answer cited: v2, Proposition 2.1 and Theorem 2.5,
 give exactly $\alpha\geq4$. Corollary~\ref{cor:completion} realizes all
 these values with irrational imaginary part.\\[4pt]
Theorem 2.6, Lemma 2.7 and Theorem 2.8: substitution and reciprocal roots
 & Clique identity reproved; prior proof Brown--Cameron (2020), citing
 earlier substitution theory. Irrational scales follow from $W[K_m]$.\\[4pt]
Theorem 2.12: induced connected supergraphs
 & Prior proof cited for $\pm\ii/m$; Corollary~\ref{cor:completion}
 extends to all $\pm\ii/\sqrt n$.\\[4pt]
Proposition 3.1 and Corollary 3.2: tree transfer and infinitely many trees
 & Prior proof Bencs (2017 preprint; 2018 journal), then Brown--Cameron
 (2020). Corollary~\ref{cor:trees} extends the rational-root consequence.\\[4pt]
Question 1, pp.~10--11: minimum graph order
 & Not covered. The source gives a 14-vertex example with roots $\pm\ii$;
 $W$ is not a minimum-order assertion.\\[4pt]
Conclusion, p.~10, and Cameron thesis, Section 3.4, p.~46:
 possible universal-vertex restriction
 & Not asserted in the source; fails for the trees of
 Corollary~\ref{cor:trees}. The case $n=1$ is already a consequence of
 Brown--Cameron's Corollary 3.2; nonsquare $n$ gives irrational examples.\\[4pt]
Cameron thesis, Conjecture 6.2.4, p.~91: imaginary parts restricted to $\pm1$
 & Refuted-and-corrected within $b^2\in\mathbb Q$ by
 Theorem~\ref{thm:classification}. Cameron's 2019 substitution and
 $\pm\ii$ examples imply rational counterexamples; Brown--Cameron
 (2020), Theorem 2.8, states them explicitly.\\[4pt]
Cameron thesis, Section 3.4, pp.~43--46: even/odd test, integer-root
 obstruction and palindromic examples
 & Prior even/odd and rational-root arguments are used in
 Theorem~\ref{thm:classification}; the palindromic construction is not covered.\\[4pt]
Imaginary parts with irrational square
 & Not covered. The rational-root argument and the classification above
 make no assertion about this class.\\
Cameron's graph-polynomial questions, fourth item: $|b|>1$
 & Not covered in general. Theorem~\ref{thm:classification} rules out the
 rational-square subclass; irrational squares remain outside its scope.\\
\end{longtable}
\normalsize

\paragraph{Verification scope.}
The classification, connected completions, induced embeddings and tree
consequence have written proofs above. The pinned Lean
file\footnote{\href{https://github.com/the-omega-institute/trureturing/blob/1f6fb1ff47df504e8dce0cd5cda2403a7aaa622e/D5/S3/StatisticalMechanics/HardCore/IrrationalIndependenceRoot.lean}{\texttt{IrrationalIndependenceRoot.lean}},
commit \texttt{1f6fb1ff47df504e8dce0cd5cda2403a7aaa622e}.}
contains \texttt{native\_polynomial}, which proves the polynomial
factorization, degree four and all independent-set counts for $W$, and
\texttt{result}, which proves existence of a finite graph with an irrational
imaginary independence root. Both are in namespace
\texttt{D5.S3.StatisticalMechanics.HardCore.IrrationalIndependenceRoot}.
They concern $W$, rather than the all-$n$ constructions or the tree consequence.
The accompanying \texttt{verify.py} enumerates every independent subset
of each graph listed in \texttt{verification.json}, checking its coefficients,
connectivity and exact root relations. Those finite checks include
$G_n$ for $1\leq n\leq4$, $J_n$ for $1\leq n\leq3$, connected completions
and induced embeddings; their precise ranges are given in that file.

\begin{thebibliography}{9}
\bibitem{BHN}
J.~I. Brown, C.~A. Hickman and R.~J. Nowakowski,
\emph{On the location of roots of independence polynomials},
Journal of Algebraic Combinatorics \textbf{19} (2004), 273--282.
\href{https://doi.org/10.1023/B:JACO.0000030703.39946.70}{doi:10.1023/B:JACO.0000030703.39946.70}.

\bibitem{BC}
J.~I. Brown and B.~Cameron,
\emph{A note on purely imaginary independence roots},
Discrete Mathematics \textbf{343} (2020), no.~12, 112113.
\href{https://doi.org/10.1016/j.disc.2020.112113}{doi:10.1016/j.disc.2020.112113}.
Posing text: \href{https://arxiv.org/abs/2001.08881v2}{arXiv:2001.08881v2}.

\bibitem{Bencs}
F.~Bencs, \emph{On trees with real-rooted independence polynomial},
Discrete Mathematics \textbf{341} (2018), no.~12, 3321--3330.
\href{https://doi.org/10.1016/j.disc.2018.06.033}{doi:10.1016/j.disc.2018.06.033}.
The tree-transfer theorem is Proposition 2.7(1) of
\href{https://arxiv.org/abs/1703.05409v1}{arXiv:1703.05409v1}.

\bibitem{Cameron}
B.~R.~K. Cameron, \emph{On the roots of independence polynomials},
Ph.D. thesis, Dalhousie University, March 2019.
\href{https://dalspace.library.dal.ca/server/api/core/bitstreams/c903bf10-142c-4ffe-965a-1927ad52197a/content}{Dalhousie institutional copy},
Conjecture 6.2.4, p.~91.

\bibitem{CameronQ}
B.~Cameron, \emph{Graph polynomials}, institutional research questions,
fourth question, undated.
\href{https://www.smcs.upei.ca/~brcameron/graphpolys.html}{University of Prince Edward Island page},
accessed November 19, 2026.
\end{thebibliography}
\smallskip
\noindent\small Licensed under Creative Commons Attribution 4.0 International
(CC BY 4.0).
\end{document}
